\section{精确反馈下的损伤保证}
\label{sec:exact}
\subsection{损伤分解与下一事件统计机制}
给定$\mathcal G_t^\pm$，真实状态仍具有高斯条件分布，条件均值记${\bm{\mu}}_t^\pm$，
协方差为${\bm{P}}_t^\pm$。发送并给定真实测量后，事件似然因子在后验归一化时消去；
静默时，积分量测噪声得到以实际预测为中心的高斯似然。该结构不要求载荷修改与当前量测独立。
令${\bm{d}}_t={\bm{\mu}}_t^--{\bm{m}}_t^-$、${\bm{D}}_t={\bm{\mu}}_t^+-{\bm{m}}_t^+$，则
\begin{align}
{\bm{D}}_t&=({\bm{I}}-{\bm{K}}_{\gamma_t}({\bm{P}}_t^-)){\bm{d}}_t-\gamma_t{\bm{K}}_1({\bm{P}}_t^-){\bm{a}}_t,\qquad
{\bm{d}}_{t+1}={\bm{A}}{\bm{D}}_t,\label{eq:bias}\\
J_k&=\E_\pi\norm{{\bm{D}}_k}^2+
\E_\pi\tr {\bm{P}}_k^+-\E_0\tr {\bm{P}}_k^+.\label{eq:split}
\end{align}
式\eqref{eq:split}由条件零均值残差的正交性得到。最后两项不能删除；
事件概率改变可能使其为负。

固定${\bm{P}}={\bm{P}}_k^+$，记下一静默位$Y=\ind_{\{\gamma_{k+1}=0\}}$，定义
\begin{align}
{\bm{\Sigma}}({\bm{P}})&={\bm{A}}{\bm{P}}{\bm{A}}^\top+{\bm{Q}}+{\bm{R}},\qquad
s({\bm{P}})=\det({\bm{I}}+{\bm{Z}}^{1/2}{\bm{\Sigma}}({\bm{P}}){\bm{Z}}^{1/2})^{-1/2},\nonumber\\
{\bm{H}}({\bm{P}})&=\tfrac12 {\bm{A}}^\top[{\bm{Z}}^{-1}+{\bm{\Sigma}}({\bm{P}})]^{-1}{\bm{A}},\qquad
W=Y/s({\bm{P}}).
\end{align}
真实条件创新的均值为${\bm{A}}{\bm{D}}_k$。高斯积分给
\begin{equation}
\E_\pi[Y\mid\mathcal G_k^+]=s({\bm{P}})e^{-{\bm{D}}_k^\top {\bm{H}}({\bm{P}}){\bm{D}}_k},
\qquad \E_0W=1,\quad \E_\pi W=\E_\pi e^{-{\bm{D}}_k^\top {\bm{H}}({\bm{P}}){\bm{D}}_k}.
\label{eq:witness}
\end{equation}
未来载荷修改发生在未来事件之后，故不改变这一关系。

式\eqref{eq:witness}是可信事件转化为损伤信息的关键。给定历史时，下一静默概率相对健康基准$s({\bm{P}})$乘以$\exp[-{\bm{D}}_k^\top {\bm{H}}({\bm{P}}){\bm{D}}_k]$；由于${\bm{A}}$可逆且${\bm{Z}}^{-1}+{\bm{\Sigma}}({\bm{P}})\succ0$，有${\bm{H}}({\bm{P}})\succ0$，每个非零偏差都会严格降低该归一化概率。审计者无需直接观察${\bm{D}}_k$，而是比较记录中${\bm{P}}_k^+$与下一事件的联合分布。

协方差标签的作用是校准事件历史对应的基准概率$s({\bm{P}})$。若只比较未校准的静默边缘概率，攻击改变不同事件历史的权重也会改变平均静默率，不能将其全部归因于隐藏偏差。式\eqref{eq:split}中的协方差分布变化项同样必须保留；定理对$R_{\rm bits}$和$R_{\rm full}$成立，是因为这些记录也保留了${\bm{P}}_k^+$所依赖的事件历史，而不是因为这些字段具有相同的信息量。

有限$k$下，P及增益取有限值。由式\eqref{eq:bias}可得
$\norm{{\bm{D}}_k}\le r_kb$。固定$b_{\rm ref}>0$，取
\[
h_*=\min_{\bm{P}}\lambda_{\min}({\bm{H}}({\bm{P}}))>0,\quad
B_*=r_kb_{\rm ref},\quad q=\frac{1-e^{-h_*B_*^2}}{B_*^2}>0.
\]
其中${\bm{A}}$可逆保证$h_*>0$。割线不等式给出
\begin{proposition}[可信下一事件给出的总体损伤上界]
对全部允许策略及$0\le b\le b_{\rm ref}$，
\begin{equation}
J_k\le \frac{1-\E_\pi W}{q}
+\E_\pi\tr {\bm{P}}_k^+-\E_0\tr {\bm{P}}_k^+.
\label{eq:certificate}
\end{equation}
\end{proposition}
\begin{proof}
在$\norm{{\bm{D}}_k}\le B_*$上，
$1-e^{-{\bm{D}}_k^\top {\bm{H}}({\bm{P}}){\bm{D}}_k}\ge q\norm{{\bm{D}}_k}^2$。
对式\eqref{eq:witness}取期望并代入式\eqref{eq:split}即得。
\end{proof}
右侧是由可信事件与已知模型标签定义的总体统计量，不需要审计者知道攻击向量或真实状态。
有限样本估计仍有统计余量；一个事件或一条轨迹不能直接认证总体损伤。

\subsection{双预算阶数及完整日志下界}
记$\Delta_m=\Delta_{R_m}$，
令$M=\max_{\bm{P}}1/s({\bm{P}})$，$\Omega_k=\max_{\bm{P}}\tr {\bm{P}}-\min_{\bm{P}}\tr {\bm{P}}$。
因$0\le W\le M$且${\bm{P}}$在记录中，标准TV期望界给
\[
J_k\le(M/q+\Omega_k)\Delta_m.
\]
另一方面，式\eqref{eq:bias}控制偏差平方；每个历史下静默概率相对健康变化为$O(b^2)$。
对有限事件前缀耦合并控制P标签的期望差，得到$J_k\le C_{\rm peak}b^2$。
该第二步不可仅从偏差半径推出而忽略协方差项。

\begin{theorem}[精确反馈的有限时域损伤阶数]
\label{thm:exact}
固定模型、有限维数、目标$k$及$b_{\rm ref}$，存在正有限常数$c,C$，
对$0<b\le b_{\rm ref}$、$0\le\varepsilon\le1$及
$R\in\{R_{\rm full},R_{\rm bits},R_m\}$，
\begin{equation}
c\min\{b^2,\varepsilon\}\le V_R^{(k)}(b,\varepsilon)
\le C\min\{b^2,\varepsilon\}.
\label{eq:exact-bound}
\end{equation}
特别地，$V_R^{(k)}(b,0)=0$。
\end{theorem}
\begin{proof}
上界为两种界的组合及记录的数据处理。
下界在目标前不修改，仅在目标真实发送后使用独立等概率符号$\sigma$和固定单位方向${\bm{v}}$，$\norm {\bm{v}}=1$，
取${\bm{a}}_k=\sigma\eta {\bm{v}}$，其他时刻为零。目标损伤为$L_{\bm{v}}\eta^2$，$L_{\bm{v}}>0$。
给定健康前史，将当前连续载荷和下一静默原子/连续载荷写成一个联合核。
该核对位移二阶连续$L^1$可微，正负混合的一阶项消去，完整日志TV不超过$C_1\eta^2$。
取$\eta=\min\{b,\sqrt{\varepsilon/C_1}\}$完成下界，联合核细节见附录\ref{app:single}。
\end{proof}
单拍构造足够支持阶数，并不声称是精确最坏策略。
常数依赖模型、$k$、$b_{\rm ref}$；$\varepsilon$与$b^2$的量纲通过模型相关常数联系，
两预算切换点不是普遍的$\varepsilon=b^2$。
三个记录的最坏损伤具有相同阶数，不表示它们的最坏值或上界常数相等。
由于$R_m$、$R_{\rm bits}$均为$R_{\rm full}$的确定性投影，数据处理给出
$V_{R_{\rm full}}^{(k)}\le V_{R_{\rm bits}}^{(k)}\le V_{R_m}^{(k)}$；
定理说明的是这三个量在固定模型和时域下具有匹配阶数。

\subsection{观察时域与信息限制}
式\eqref{eq:witness}使用下一拍可信事件以及目标协方差标签。
因此，定理\ref{thm:exact}不直接适用于只观察到目标拍的日志，或删去历史校准标签后的单个事件边缘。
这一充分保证不声称已确定所有审计方式的最小记录。
精确协方差标签由事件历史决定，但并不天然构成低比特压缩；
本文未给出其固定字长编码或通信成本最优性。
